\section{Ausgangssituation}

Mitarbeitende der Primetals Technologies Austria GmbH müssen derzeit mehrere
Systeme verwenden, um Informationen über die IT-Infrastruktur des Unternehmens
zu erhalten und potenzielle Probleme zu identifizieren. Nach Angaben der
Mitarbeitenden sind relevante Informationen auf verschiedene Quellen verteilt,
vor allem auf \gls{mssql}-Datenbanken und Weboberflächen, wodurch das Monitoring
der Infrastruktur zeitaufwendig und ineffizient wird. Mitarbeitende, die bereits
mit diesen Systemen gearbeitet haben, berichteten außerdem, dass wichtige
Informationen aufgrund des fehlenden zentralen Überblicks leicht übersehen werden
können.

\subsection{Unternehmen}

Die Primetals Technologies Austria GmbH ist ein internationales
Ingenieurunternehmen, das sich auf Technologien, Anlagen und automatisierte
Lösungen für die Metallindustrie spezialisiert hat. Das Unternehmen entwickelt
Lösungen für die Stahlproduktion mit Fokus auf Produktivitätssteigerung,
Nachhaltigkeit und Digitalisierung~\cite{primetals2026}.

\subsection{Infrastrukturlandschaft}

Die Infrastruktur der Primetals Technologies Austria GmbH besteht aus mehreren
Quellen, die Informationen über unterschiedliche Unternehmensressourcen
enthalten. Die zentralen \gls{mssql}-Datenbanken speichern strukturierte
Informationen über Geräte, Standorte und Netzwerke in Form von Tabellen und
Views. Darüber hinaus stellen Webanwendungen Informationen über Geräte,
Azure-Dienste und den Verbindungsstatus verschiedener Netzwerke bereit.

\section{Problemstellung}

Die aktuelle Infrastruktur erfordert, dass Mitarbeitende Daten manuell aus
unabhängigen Quellen zusammenführen, um sich ein vollständiges Bild der
Unternehmensinfrastruktur zu machen. Dieser manuelle Prozess kostet Zeit und
erhöht das Risiko von Fehlern bei der Datenanalyse. Dadurch fehlt eine
zuverlässige Möglichkeit, den aktuellen Zustand der IT-Infrastruktur des
Unternehmens effizient zu überwachen.

\section{Konzept}

Das vorgestellte Konzept basiert auf einer zentralen Anwendung, die
Informationen aus bestehenden Quellen sammelt und aggregiert. Die erfassten
Daten werden über eine zentrale Datasource-Schicht verarbeitet, die eine
einheitliche Schnittstelle für den Zugriff auf Infrastrukturdaten bereitstellt.
Statt mehrere Systeme manuell zu durchsuchen, erhalten Mitarbeitende eine
zentrale Visualisierung, die sie bei der Identifikation von Problemen in der
technischen Landschaft des Unternehmens unterstützt.

\section{Vorgeschlagene Lösung}

Ziel ist ein \gls{grafana}-Plugin, das das Global Infrastructure Overview
(GIO)-System implementiert. Das Plugin ist in mehrere Schichten unterteilt, um
Datenverarbeitung, Kommunikation und Visualisierung voneinander zu trennen.

Die Common-Schicht enthält Repositories und Provider, die die erfassten Daten
verarbeiten und aggregieren, bevor sie über die \gls{grafana}-Go-API-Schnittstelle
verfügbar gemacht werden.

Die Datasource-Schicht kann erweitert werden, um Daten aus \gls{mssql}-Datenbanken
und APIs abzurufen. Die gesammelten Informationen werden in eine einheitliche
Struktur überführt und an die Visualisierungsschicht weitergeleitet, die auf die
Daten unabhängig von der ursprünglichen Datenstruktur zugreifen kann.

Die Visualisierungsschicht stellt die Daten der Primetals Technologies Austria
GmbH auf einer Weltkarte sowie in Detailansichten für weiterführende
Informationen dar. Die Ansicht verändert sich abhängig von Zoom- und
Klickaktionen. Zusätzlich helfen Filter- und Suchfunktionen den Mitarbeitenden
dabei, ein Problem zu lokalisieren und zu analysieren.

Der Zugriff auf \gls{grafana} ist durch die vorhandenen Authentifizierungs- und
Autorisierungsmechanismen abgesichert, sodass nur berechtigte Benutzer auf das
Plugin zugreifen können.

\section{Aufbau der Arbeit}

Die Arbeit ist in mehrere Kapitel gegliedert. Nach der Einleitung werden die
technischen Grundlagen von \gls{grafana} und den erforderlichen Technologien
vorgestellt. Das folgende Kapitel beschreibt die Architektur und die
Implementierung der entwickelten Lösung. Abschließend enthält die Arbeit eine
Evaluierung des Systems, eine Diskussion der Ergebnisse sowie mögliche
Erweiterungen für die Zukunft.